(1) 在单激发子空间 \(\{|1,0_E\rangle,|0,1_E\rangle\}\) 中，可取
\[
 U_{\rm 1exc}=\begin{pmatrix}\sqrt\eta&-\sqrt{1-\eta}\\
 \sqrt{1-\eta}&\sqrt\eta\end{pmatrix},\qquad
 U_{\rm 1exc}^\dagger U_{\rm 1exc}=I .
\]
结合真空态不变，令环境测量基为 \(|0_E\rangle,|1_E\rangle\)，得到
\[
 K_0=\langle0_E|U|0_E\rangle=
 \begin{pmatrix}1&0\\0&\sqrt\eta\end{pmatrix},\qquad
 K_1=\langle1_E|U|0_E\rangle=
 \begin{pmatrix}0&\sqrt{1-\eta}\\0&0\end{pmatrix}.
\]
因为 \(K_0^\dagger K_0+K_1^\dagger K_1=I\)，映射
\(\Lambda_\eta(\rho)=\sum_{j=0}^1K_j\rho K_j^\dagger\)
完全正且保迹。设
\(\rho=\begin{pmatrix}1-p&z\\z^*&p\end{pmatrix}\)，
其中 \(0\le p\le1\)、\(|z|^2\le p(1-p)\)，则
\[
 \boxed{\Lambda_\eta(\rho)=
 \begin{pmatrix}1-\eta p&\sqrt\eta\,z\\
 \sqrt\eta\,z^*&\eta p\end{pmatrix}}.
\]
若 \(\rho=(I+x\sigma_x+y\sigma_y+z_B\sigma_z)/2\)，且
\(\sigma_z=|0\rangle\langle0|-|1\rangle\langle1|\)，则
\((x,y,z_B)\mapsto(\sqrt\eta x,\sqrt\eta y,1-\eta+\eta z_B)\)。
该通道非幺正、非幺正混合：最大混合态不保持不变。

(2) 记 \(\sigma_-=|0\rangle\langle1|\)、\(\sigma_+=\sigma_-^\dagger\)。
对(1)式逐元微分，使用 \(\dot\eta=-\Gamma\eta\)，在相互作用绘景得到
\[
 \boxed{\dot\rho=\Gamma\left(\sigma_-\rho\sigma_+
 -\tfrac12\{\sigma_+\sigma_-,\rho\}\right)}.
\]
因而 \(\dot p=-\Gamma p\)、\(\dot z=-\Gamma z/2\)，
\(p(t)=p(0)e^{-\Gamma t}\)，
\(z(t)=z(0)e^{-\Gamma t/2}\)；若返回薛定谔绘景，非对角元另有 \(e^{i\omega t}\) 或其共轭的自由进动相位，具体符号取决于矩阵元的排列。
直接将(1)式连续施加两次得
\(\Lambda_{\eta_2}\circ\Lambda_{\eta_1}=\Lambda_{\eta_1\eta_2}\)，
故 \(\Lambda_{t+s}=\Lambda_t\circ\Lambda_s\)。
系统平均能量为 \(E(t)=\hbar\omega p(t)\)，
\[
 \boxed{\frac{dE}{dt}=-\Gamma\hbar\omega p(t)\le0}.
\]
损失的能量流入环境；系统能量不单独守恒，联合系统在旋波近似下守恒激发能。

(3) 取基底顺序 \((|00\rangle,|01\rangle,|10\rangle,|11\rangle)\)。
通道在 \(B\) 上作用后
\[
 \rho_{AB}=\frac12
 \begin{pmatrix}
 1&0&0&\sqrt\eta\\
 0&0&0&0\\
 0&0&1-\eta&0\\
 \sqrt\eta&0&0&\eta
 \end{pmatrix}.
\]
对第二个比特作部分转置，只需将矩阵元
\(\langle ij|\rho|kl\rangle\) 的 \(j,l\) 互换：
\[
 \rho_{AB}^{T_B}=\frac12
 \begin{pmatrix}
 1&0&0&0\\
 0&0&\sqrt\eta&0\\
 0&\sqrt\eta&1-\eta&0\\
 0&0&0&\eta
 \end{pmatrix}.
\]
中间 \(2\times2\) 块的本征值为 \(1/2,-\eta/2\)，
因此全谱为
\[
 \boxed{\operatorname{spec}(\rho^{T_B})=
 \{1/2,1/2,\eta/2,-\eta/2\}},\qquad
 \boxed{\mathcal N=\eta/2}.
\]
任意有限 \(t\) 有 \(\eta>0\)，故部分转置非正，双比特仍纠缠；仅在 \(t\to\infty\) 时负性趋零，不存在有限时间的纠缠突然死亡。注意此结论针对给定 Bell 初态及单边零温振幅阻尼，并非所有初态的普遍定理。

(4) 记 \(p=|\beta|^2\)，由(1)式直接求
\[
 F_\psi=\langle\psi|\Lambda_\eta(|\psi\rangle\langle\psi|)|\psi\rangle
 =(1-p)(1-\eta p)+\eta p^2+2\sqrt\eta\,p(1-p).
\]
对单比特 Haar 随机纯态，\(p\) 在 \([0,1]\) 上均匀分布，故
\[
 \boxed{F_{\rm av}=\int_0^1F_\psi\,dp
 =\frac{3+2\sqrt\eta+\eta}{6}}.
\]
独立地，最大混合态的纠缠保真度可用 Kraus 迹公式
\[
 F_e=\sum_j|\operatorname{Tr}[(I/2)K_j]|^2
 =\frac{|\operatorname{Tr}K_0|^2+|\operatorname{Tr}K_1|^2}{4}
 =\frac{(1+\sqrt\eta)^2}{4}.
\]
于是 \((2F_e+1)/3=(3+2\sqrt\eta+\eta)/6\)，
与直接积分完全一致。\(t=0\) 时 \(\eta=1\)、\(F_{\rm av}=1\)；
\(t\to\infty\) 时输出恒为 \(|0\rangle\langle0|\)，
\(F_{\rm av}=1/2\)，而非错误地认为完全衰减仍保留未知量子态。
所有概率、负性与保真度均无量纲；\(\Gamma^{-1}\) 为时间量纲。
